%%%%%%%%%%%%%%%%%%%%%%%%%%%%%%%%%%%%%%%
% Assurance- and ETB-related terms and acronyms
%
% input this glossary early in your preamble after the glossaries package and options
%    e.g.,
%		\usepackage[toc,acronym,style=altlist]{glossaries}
%		\makeglossaries
%		\renewcommand{\glspostdescription}{}                 % suggested
%		\renewcommand{\glsnamefont}[1]{\textbf{#1}}      % suggested
%		%%%%%%%%%%%%%%%%%%%%%%%%%%%%%%%%%%%%%%%
% Assurance- and ETB-related terms and acronyms
%
% input this glossary early in your preamble after the glossaries package and options
%    e.g.,
%		\usepackage[toc,acronym,style=altlist]{glossaries}
%		\makeglossaries
%		\renewcommand{\glspostdescription}{}                 % suggested
%		\renewcommand{\glsnamefont}[1]{\textbf{#1}}      % suggested
%		%%%%%%%%%%%%%%%%%%%%%%%%%%%%%%%%%%%%%%%
% Assurance- and ETB-related terms and acronyms
%
% input this glossary early in your preamble after the glossaries package and options
%    e.g.,
%		\usepackage[toc,acronym,style=altlist]{glossaries}
%		\makeglossaries
%		\renewcommand{\glspostdescription}{}                 % suggested
%		\renewcommand{\glsnamefont}[1]{\textbf{#1}}      % suggested
%		%%%%%%%%%%%%%%%%%%%%%%%%%%%%%%%%%%%%%%%
% Assurance- and ETB-related terms and acronyms
%
% input this glossary early in your preamble after the glossaries package and options
%    e.g.,
%		\usepackage[toc,acronym,style=altlist]{glossaries}
%		\makeglossaries
%		\renewcommand{\glspostdescription}{}                 % suggested
%		\renewcommand{\glsnamefont}[1]{\textbf{#1}}      % suggested
%		\input{../glossary-Assurance}
%
% Don't forget to add e.g. ../../glossary-Assurance dependency to your Makefile if you
% are adding terms to a glossary while working on your document.

\newglossaryentry{analysis tool}{
	name={analysis tool},
	description={a tool that is used to perform some analysis or test the result of which may serve as evidence in support of an AC}
}

\newglossaryentry{argument pattern}{
	name={argument pattern},
	description={a parameterized template described in the GSN standard for an AC fragment possibly involving structural abstraction}
}

\newglossaryentry{assurance case}{
	name={assurance case},
	description={a documented body of evidence that provides a convincing and valid argument that a system is adequately dependable for a given application in a given environment; provides a set of auditable claims, arguments, and evidence created to support the claim that a defined system/service will satisfy particular requirements}
}

\newglossaryentry{assurance deficit}{
	name={assurance deficit},
	description={a documented body of evidence that provides a convincing and valid argument that a system is adequately dependable for a given application in a given environment; provides a set of auditable claims, arguments, and evidence created to support the claim that a defined system/service will satisfy particular requirements}
}

\newglossaryentry{assurance case argument pattern}{
        name={assurance case argument pattern},
        description={a parameterized expression representing logical relationships among goals (or claims) and evidence, expressed in a form of GSN with extensions for structural abstraction (see also: argument pattern}
}

\newglossaryentry{batch mode}{
        name={batch mode (O-ETB)},
        description={a mode of operation of O-ETB that is activated when the \texttt{-{}-command} or \texttt{-c} option is used to
			execute a single command or to run pre-defined command procedure or script (by executing the \texttt{proc} or \texttt{script}
			as the single command);
			other options that establish the context for the execution of the command are performed before the command;
			O-ETB exits after performing the specified command \cite{O-ETB-UM}}
}

\newglossaryentry{CAP artefact}{
	name={CAP artefact (O-ETB)},
	description={an object consisting of one or more files located in the CAP directory or a project directory within the CAP that has been produced to fulfill a requirement for certification evaluation}
}

\newglossaryentry{CAP Renderer}{
	name={CAP Renderer (O-ETB)},
	description={a component of O-ETB that renders information in the CASES Repository and the EVIDENCE Repository in an external form suitable for review as part of a certification evaluation}
}

\newglossaryentry{CASES Repository}{
	name={CASES Repository (O-ETB)},
	description={a persistent store created by O-ETB that contains modular ACs instantiated from one or more argument patterns and referring to evidence records in the ER \cite{O-ETB-UM}}
}

\newglossaryentry{certification assurance package}{
        name={certification assurance package (CAP) (O-ETB)},
        description={a collection of documents (textual and/or graphical and electronically searchable),
			including those generated by the O-ETB,
			required for evaluation that present required evidence in a form acceptable to a certification authority \cite{O-ETB-UM}}
}

\newglossaryentry{certification project}{
        name={certification project},
        description={a subdivision within the CAP for sets of artifacts corresponding to different certification activities \cite{O-ETB-UM}}
}

\newglossaryentry{certification requirements}{
	name={certification requirements},
	description={a set of requirements prescribing methods and/or the content and presentation of evidence needed to justify granting certification according to some particular domain or standard}
}

\newglossaryentry{claims-argument-evidence}{
        name={claims-argument-evidence (CAE)},
        description={a structured approach which aids the visibility, review and assessment of safety case documentation \cite{CAE}}
}

\newglossaryentry{command mode}{
        name={command mode (O-ETB)},
        description={a command-based interface that can be used interactively and/or in a batch mode to execute single commands or to run pre-defined command scripts \cite{O-ETB-UM}}
}

\newglossaryentry{command procedure}{
        name={command procedure (O-ETB)},
        description={a pre-defined named sequence of O-ETB commands, defined along with other such procedures
			in the \texttt{procs\_etb.pl} file,
			that can be invoked by using the \texttt{proc} command either in the interactive command interpreter
			or in a \texttt{-{}-command} option in batch mode \cite{O-ETB-UM}}
}

\newglossaryentry{command script}{
        name={command script (O-ETB)},
        description={a user-defined sequence of O-ETB commands contained in a named file,
			that can be invoked by using the \texttt{script} command either in the interactive command interpreter
			or in a \texttt{-{}-command} option in batch mode \cite{O-ETB-UM}}
}

\newglossaryentry{common attack pattern enumeration and classification}{
        name={common attack pattern enumeration and classification (CAPEC)},
        description={a comprehensive dictionary of known patterns of attack employed by adversaries to exploit known weaknesses in cyber-enabled capabilities \cite{CAPEC}}
}

\newglossaryentry{common criteria}{
        name={common criteria (CC)},
        description={Common Criteria for Information Security Technology Evaluation \cite{ISOstd-15408-1,ISOstd-15408-2,ISOstd-15408-3}}
}

\newglossaryentry{common vulnerabilities enumeration}{
        name={common vulnerabilities enumeration (CVE)},
        description={catalog of publicly disclosed cybersecurity vulnerabilities \cite{CVE}}
}

\newglossaryentry{common weaknesses enumeration}{
        name={common weaknesses enumeration (CWE)},
        description={the most severe and prevalent weaknesses behind the CVE records \cite{CWE}}
}

\newglossaryentry{correct by construction}{
	name={correct by construction (CxC)},
	description={rigorous techniques applied to the construction of some artifact or product that logically ensure that the item correctly adheres to a specified property or properties}
}

\newglossaryentry{deductive verification}{
	name={deductive verification},
	description={deductive software verification aims at formally verifying that all possible behaviors of a given program satisfy formally defined, possibly complex properties, where the verification process is based on some form of logical inference, i.e., deduction}
}

\newglossaryentry{ETB Core Workflow Engine}{
	name={ETB Core Workflow Engine (O-ETB)},
	description={a component of O-ETB that is the top-level driver for the instantiation of assurance case argument patterns and the passing of actual parameter values \cite{O-ETB-UM}}
}

\newglossaryentry{ETB Developer/Administrator}{
	name={ETB Developer/Administrator (O-ETB)},
	description={the role of a person that develops ETB implementation or administers an ETB implementation}
}

\newglossaryentry{evaluation assurance level}{
        name={evaluation assurance level (EAL) (CC)},
        description={set of assurance requirements that represent a point on the Common Criteria predefined assurance scale (EALs 1..7) \cite{CC,ISOstd-15408-1}}
}

\newglossaryentry{evaluator (certifier)}{
	name={evaluator (certifier)},
	description={a qualified expert who examines evidence presented in support of an application by a developer to certify a product}
}

\newglossaryentry{evidence}{
	name={evidence},
	description={concrete empirical data that may be used to support arguments for claims}
}

\newglossaryentry{evidence category}{
	name={evidence category (O-ETB)},
	description={an identifier for a set of values that a particular O-ETB evidence item may represent and which is produced by a particular method or tool \cite{O-ETB-UM}}
}

\newglossaryentry{EVIDENCE Repository}{
	name={EVIDENCE Repository (O-ETB)},
	description={a persistent store created by O-ETB and updated by tool agents that contains evidence records that represent placeholders for evidence or the concrete result of an assurance activity or the output of an evidential tool \cite{O-ETB-UM}}
}

\newglossaryentry{evidence specifier}{
	name={evidence specifier (O-ETB)},
	description={unique identifier for an instance of an evidence category occurring in the EVIDENCE Repository, see D2.2 p26 \cite{O-ETB-UM}}
}

\newglossaryentry{evidential tool bus}{
	name={evidential tool bus (ETB)},
	description={a system such as O-ETB that constructs an assurance case and supporting evidence according to a collection of argument patterns and invokes and records the results from external evidential tools \cite{Rushby05}}
}

\newglossaryentry{Frama-C}{
	name={Frama-C},
	description={Framework for Modular Analysis of C programs;
		an open-source extensible and collaborative platform dedicated to source-code analysis of C software}
}

\newglossaryentry{Frama-PLC}{
	name={Frama-PLC},
	description={performs syntactical and semantic checks of PLC application software}
}

\newglossaryentry{goal structuring notation}{
        name={goal structuring notation (GSN)},
        description={a graphical argumentation notation that explicitly represents the individual elements of any safety argument (requirements, claims, evidence and context) and the relationships that exist between these elements \cite{kelly2004goal}}
}

\newglossaryentry{known answer test}{
	name={known answer test (KAT)},
	description={a test of implementation correctness that employs pairs of parameter values and results that are calculated or otherwise known independently of the implementation being tested}
}

\newglossaryentry{model checking}{
	name={model checking},
	description={a form of automatic formal verification that employs a model checker tool to verify temporal behavior of a system
	as modeled by a symbolic state machine (transition system);
	a model checker may be used as an evidential tool to generate evidence items in an assurance case by O-ETB}
}

\newglossaryentry{O-ETB command (O-ETB)}{
	name={O-ETB command (O-ETB)},
	description={a command defined and accepted by the O-ETB command interpreter, either interactively or in a command procedure or script, to invoke an O-ETB function}
}

\newglossaryentry{Open Evidential Tool Bus}{
        name={Open Evidential Tool Bus (O-ETB)},
        description={an implementation of the ETB concept \cite{Rushby05} derived from AM-ETB \cite{CITADEL-D5.1} and adapted and extended for IoMT and MedSecurance by The Open Group \cite{O-ETB-UM}}
}

\newglossaryentry{privacy}{
	name={privacy},
	description={right of individuals to control or influence what information related to them may be collected and stored and by whom and to whom that information may be disclosed (source: ISO  TS 17574:2009)}
}

\newglossaryentry{protection profile}{
	name={protection profile (CC)},
	description={a requirements statement put together using CC constructs. Protection Profiles are generally published by governments for a specific technology type \cite{CC,ISOstd-15408-1}}
}

\newglossaryentry{reliability}{
	name={reliability},
	description={ability of a system or component to perform its required functions under stated conditions for a specified period of time (source: ISO/IEC 27040:2015)}
}

\newglossaryentry{REPOSITORY}{
        name={REPOSITORY (O-ETB)},
        description={a directory containing the persistent CASES Repository and the EVIDENCE Repository created and used by O-ETB \cite{O-ETB-UM}}
}

\newglossaryentry{resilience}{
	name={resilience},
	description={the condition of the system being able to avoid, absorb and/or manage dynamic adversarial conditions while completing assigned mission(s), and to reconstitute operational capabilities after casualties (source: IIC)}
}

\newglossaryentry{RESTful API}{
	name={RESTful API (or REST API)},
	description={an interface based on the REST architectural style that allows applications to communicate over the Internet}
}

\newglossaryentry{safety}{
	name={safety},
	description={the condition of the system operating without causing unacceptable risk of physical injury or damage to the health of people, either directly, or indirectly as a result of damage to property or to the environment (source: ISO/IEC Guide 55:1999)}
}

\newglossaryentry{security}{
	name={security},
	description={condition of the system being protected from unintended or unauthorized access, change or destruction. [Note: Security is a property of a system by which confidentiality, integrity, availability, accountability, authenticity, and reliability are achieved (ISO TR 15443-1:2012)}
}

\newglossaryentry{security target}{
	name={security target (CC)},
	description={implementation-dependent statement of security needs for a specific identified Target of Evaluation (TOE) \cite{CC,ISOstd-15408-1}}
}

\newglossaryentry{source repository}{
        name={source repository (O-ETB)},
        description={a directory structure containing the development directories of source code, documentation, build and test support for the O-ETB,
			and a template for the persistent database directories (KB, REPOSITORY and CAP) used and created by O-ETB \cite{O-ETB-UM}}
}

\newglossaryentry{stress test}{
	name={stress test},
	description={a test of a program or system that demands that it operate at high capacity, in volume and/or time, without failure}
}

\newglossaryentry{structured assurance case metamodel}{
	name={structured assurance case metamodel (SACM)},
	description={OMG standard metamodel for representing structured assurance cases and a graphical notation for depicting an assurance case \cite{SACM,SACM2}}                          
}

\newglossaryentry{system model}{
	name={system model (O-ETB)},
	description={a formal model of the architecture, properties and interactions of a system's components, used by O-ETB to provide information about the system that is needed during the assurance case argument pattern instantiation process}
}

\newglossaryentry{target of evaluation}{
	name={target of evaluation (CC)},
	description={an information system, part of a system or product, and all associated documentation, that is the subject of a security evaluation,
		in particular under the CC \cite{CC,ISOstd-15408-1}}
}

\newglossaryentry{tool agent}{
	name={tool agent (O-ETB)},
	description={a program invoked in a standard way by O-ETB that is intended to manage the execution of an external evidential tool,
		or performance of another assurance activity,
		to generate or validate an evidence item by invoking an external evidential tool or other assurance activity;
		the tool agent implements an interface to the external tool that is based on the specific data needs and interactions of that tool
		the tool agent updates the item in the EVIDENCE Repository that corresponds to an instance of the evidence caegory referenced by
		one or more particular locations in an assurance case \cite{O-ETB-UM}}
}

\newglossaryentry{trustworthiness}{
	name={trustworthiness},
	description={the degree of confidence one has that the system performs as expected with characteristics including safety, security, privacy, reliability and resilience in the face of environmental disruptions, human errors, system faults and attacks \cite{IISF}}
}

\newglossaryentry{workflow}{
	name={workflow},
	description={a sequence of industrial, administrative or other processes through which a piece of work passes from initiation to completion}
}

\newglossaryentry{workflow activity}{
	name={workflow activity},
	description={a specific task in a workflow system that is to be carried out by an automated or human agent; in the context of O-ETB, an assurance activity}
}

\newglossaryentry{workflow program}{
	name={workflow program},
	description={a sequence of statements expressed in a workflow language that defines the steps of a business process that are to be carried out as workflow activities by automated or human agents;
		O-ETB executes a workflow of assurance activities and tool invocations that is implicit in the procedure by which argument patterns are instantiated}
}

\newacronym{AC}{AC}{Assurance Case}
\newacronym{ACSL}{ACSL}{ANSI C Specification Language}
\newacronym{A-G}{A-G}{Assume-Guarantee}
\newacronym{API}{API}{Application Programming Interface}
\newacronym{CAE}{CAE}{Claims-Argument-Evidence}
\newacronym{CAP}{CAP}{Certification Assurance Package \cite{O-ETB-UM}}
\newacronym{CC}{CC}{Common Criteria for Information Technology Security Evaluation \cite{ISOstd-15408-1,ISOstd-15408-2,ISOstd-15408-3}}
\newacronym{CR}{CR}{CASES Repository (O-ETB) \cite{O-ETB-UM}}
\newacronym{CVE}{CVE}{Common Vulnerabilities Enumeration \cite{CVE}}
\newacronym{CWE}{CWE}{Common Weakness Enumeration \cite{CWE}}
\newacronym{CAPEC}{CAPEC}{Common Attack Pattern Enumeration and Classification \cite{CAPEC}}
\newacronym{CxC}{CxC}{Correct by Construction}
\newacronym{EAL}{EAL}{Evaluation Assurance Level \cite{ISOstd-15408-1}}
\newacronym{ER}{ER}{EVIDENCE Repository (O-ETB) \cite{O-ETB-UM}}
\newacronym{ETB}{ETB}{Evidential Tool Bus \cite{Rushby05}}
\newacronym{FV}{FV}{Formal Verification}
\newacronym{GSN}{GSN}{Goal Structuring Notation \cite{GSN_standard,GSN_standard_3}}
\newacronym{HiLARe}{HiLARe}{High-level ACSL Requirements}
\newacronym{IIRA}{IIRA}{Industrial Internet Reference Architecture \cite{IIRA}}
\newacronym{IISF}{IISF}{Industrial Internet Security Framework \cite{IISF}}
\newacronym{IIV}{IIV}{The Industrial Internet Vocabulary \cite{IIVocab}}
\newacronym{IoMT}{IoMT}{Internet of Medical Things \cite{MedSecurance-WWW}}
\newacronym{KB}{KB}{Knowledge Base \cite{O-ETB-UM}}
\newacronym{ltl}{LTL}{Linear-time Temporal Logic}
\newacronym{O-ETB}{O-ETB}{Open Evidential Tool Bus \cite{O-ETB-UM}}
\newacronym{OMG}{OMG}{Object Management Group}
\newacronym{MedSed}{MedSec}{MedSecurance Project}
\newacronym{MS}{MS}{MedSecurance Project}
\newacronym{PP}{PP}{Protection Profile \cite{CC,ISOstd-15408-1}}
\newacronym{REST}{REST}{Representational State Transfer}
\newacronym{RMV}{RMV}{Runtime Monitoring and Verification}
\newacronym{SACM}{SACM}{Structured Assurance Case Metamodel \cite{SACM,SACM2}}
\newacronym{SAR}{SAR}{Security Assurance Requirement \cite{CC,ISOstd-15408-1,ISOstd-15408-3}}
\newacronym{ST}{ST}{Security Target \cite{CC,ISOstd-15408-1}}
\newacronym{SFR}{SFR}{Security Functional Requirement \cite{CC,ISOstd-15408-1,ISOstd-15408-2}}
\newacronym{TA}{TA}{Tool Agent \cite{O-ETB-UM}}
\newacronym{TOE}{TOE}{Target of Evaluation \cite{CC,ISOstd-15408-1}}
\newacronym{TVRA}{TVRA}{Threat Vulnerability Risk Analysis}

%
% Don't forget to add e.g. ../../glossary-Assurance dependency to your Makefile if you
% are adding terms to a glossary while working on your document.

\newglossaryentry{analysis tool}{
	name={analysis tool},
	description={a tool that is used to perform some analysis or test the result of which may serve as evidence in support of an AC}
}

\newglossaryentry{argument pattern}{
	name={argument pattern},
	description={a parameterized template described in the GSN standard for an AC fragment possibly involving structural abstraction}
}

\newglossaryentry{assurance case}{
	name={assurance case},
	description={a documented body of evidence that provides a convincing and valid argument that a system is adequately dependable for a given application in a given environment; provides a set of auditable claims, arguments, and evidence created to support the claim that a defined system/service will satisfy particular requirements}
}

\newglossaryentry{assurance deficit}{
	name={assurance deficit},
	description={a documented body of evidence that provides a convincing and valid argument that a system is adequately dependable for a given application in a given environment; provides a set of auditable claims, arguments, and evidence created to support the claim that a defined system/service will satisfy particular requirements}
}

\newglossaryentry{assurance case argument pattern}{
        name={assurance case argument pattern},
        description={a parameterized expression representing logical relationships among goals (or claims) and evidence, expressed in a form of GSN with extensions for structural abstraction (see also: argument pattern}
}

\newglossaryentry{batch mode}{
        name={batch mode (O-ETB)},
        description={a mode of operation of O-ETB that is activated when the \texttt{-{}-command} or \texttt{-c} option is used to
			execute a single command or to run pre-defined command procedure or script (by executing the \texttt{proc} or \texttt{script}
			as the single command);
			other options that establish the context for the execution of the command are performed before the command;
			O-ETB exits after performing the specified command \cite{O-ETB-UM}}
}

\newglossaryentry{CAP artefact}{
	name={CAP artefact (O-ETB)},
	description={an object consisting of one or more files located in the CAP directory or a project directory within the CAP that has been produced to fulfill a requirement for certification evaluation}
}

\newglossaryentry{CAP Renderer}{
	name={CAP Renderer (O-ETB)},
	description={a component of O-ETB that renders information in the CASES Repository and the EVIDENCE Repository in an external form suitable for review as part of a certification evaluation}
}

\newglossaryentry{CASES Repository}{
	name={CASES Repository (O-ETB)},
	description={a persistent store created by O-ETB that contains modular ACs instantiated from one or more argument patterns and referring to evidence records in the ER \cite{O-ETB-UM}}
}

\newglossaryentry{certification assurance package}{
        name={certification assurance package (CAP) (O-ETB)},
        description={a collection of documents (textual and/or graphical and electronically searchable),
			including those generated by the O-ETB,
			required for evaluation that present required evidence in a form acceptable to a certification authority \cite{O-ETB-UM}}
}

\newglossaryentry{certification project}{
        name={certification project},
        description={a subdivision within the CAP for sets of artifacts corresponding to different certification activities \cite{O-ETB-UM}}
}

\newglossaryentry{certification requirements}{
	name={certification requirements},
	description={a set of requirements prescribing methods and/or the content and presentation of evidence needed to justify granting certification according to some particular domain or standard}
}

\newglossaryentry{claims-argument-evidence}{
        name={claims-argument-evidence (CAE)},
        description={a structured approach which aids the visibility, review and assessment of safety case documentation \cite{CAE}}
}

\newglossaryentry{command mode}{
        name={command mode (O-ETB)},
        description={a command-based interface that can be used interactively and/or in a batch mode to execute single commands or to run pre-defined command scripts \cite{O-ETB-UM}}
}

\newglossaryentry{command procedure}{
        name={command procedure (O-ETB)},
        description={a pre-defined named sequence of O-ETB commands, defined along with other such procedures
			in the \texttt{procs\_etb.pl} file,
			that can be invoked by using the \texttt{proc} command either in the interactive command interpreter
			or in a \texttt{-{}-command} option in batch mode \cite{O-ETB-UM}}
}

\newglossaryentry{command script}{
        name={command script (O-ETB)},
        description={a user-defined sequence of O-ETB commands contained in a named file,
			that can be invoked by using the \texttt{script} command either in the interactive command interpreter
			or in a \texttt{-{}-command} option in batch mode \cite{O-ETB-UM}}
}

\newglossaryentry{common attack pattern enumeration and classification}{
        name={common attack pattern enumeration and classification (CAPEC)},
        description={a comprehensive dictionary of known patterns of attack employed by adversaries to exploit known weaknesses in cyber-enabled capabilities \cite{CAPEC}}
}

\newglossaryentry{common criteria}{
        name={common criteria (CC)},
        description={Common Criteria for Information Security Technology Evaluation \cite{ISOstd-15408-1,ISOstd-15408-2,ISOstd-15408-3}}
}

\newglossaryentry{common vulnerabilities enumeration}{
        name={common vulnerabilities enumeration (CVE)},
        description={catalog of publicly disclosed cybersecurity vulnerabilities \cite{CVE}}
}

\newglossaryentry{common weaknesses enumeration}{
        name={common weaknesses enumeration (CWE)},
        description={the most severe and prevalent weaknesses behind the CVE records \cite{CWE}}
}

\newglossaryentry{correct by construction}{
	name={correct by construction (CxC)},
	description={rigorous techniques applied to the construction of some artifact or product that logically ensure that the item correctly adheres to a specified property or properties}
}

\newglossaryentry{deductive verification}{
	name={deductive verification},
	description={deductive software verification aims at formally verifying that all possible behaviors of a given program satisfy formally defined, possibly complex properties, where the verification process is based on some form of logical inference, i.e., deduction}
}

\newglossaryentry{ETB Core Workflow Engine}{
	name={ETB Core Workflow Engine (O-ETB)},
	description={a component of O-ETB that is the top-level driver for the instantiation of assurance case argument patterns and the passing of actual parameter values \cite{O-ETB-UM}}
}

\newglossaryentry{ETB Developer/Administrator}{
	name={ETB Developer/Administrator (O-ETB)},
	description={the role of a person that develops ETB implementation or administers an ETB implementation}
}

\newglossaryentry{evaluation assurance level}{
        name={evaluation assurance level (EAL) (CC)},
        description={set of assurance requirements that represent a point on the Common Criteria predefined assurance scale (EALs 1..7) \cite{CC,ISOstd-15408-1}}
}

\newglossaryentry{evaluator (certifier)}{
	name={evaluator (certifier)},
	description={a qualified expert who examines evidence presented in support of an application by a developer to certify a product}
}

\newglossaryentry{evidence}{
	name={evidence},
	description={concrete empirical data that may be used to support arguments for claims}
}

\newglossaryentry{evidence category}{
	name={evidence category (O-ETB)},
	description={an identifier for a set of values that a particular O-ETB evidence item may represent and which is produced by a particular method or tool \cite{O-ETB-UM}}
}

\newglossaryentry{EVIDENCE Repository}{
	name={EVIDENCE Repository (O-ETB)},
	description={a persistent store created by O-ETB and updated by tool agents that contains evidence records that represent placeholders for evidence or the concrete result of an assurance activity or the output of an evidential tool \cite{O-ETB-UM}}
}

\newglossaryentry{evidence specifier}{
	name={evidence specifier (O-ETB)},
	description={unique identifier for an instance of an evidence category occurring in the EVIDENCE Repository, see D2.2 p26 \cite{O-ETB-UM}}
}

\newglossaryentry{evidential tool bus}{
	name={evidential tool bus (ETB)},
	description={a system such as O-ETB that constructs an assurance case and supporting evidence according to a collection of argument patterns and invokes and records the results from external evidential tools \cite{Rushby05}}
}

\newglossaryentry{Frama-C}{
	name={Frama-C},
	description={Framework for Modular Analysis of C programs;
		an open-source extensible and collaborative platform dedicated to source-code analysis of C software}
}

\newglossaryentry{Frama-PLC}{
	name={Frama-PLC},
	description={performs syntactical and semantic checks of PLC application software}
}

\newglossaryentry{goal structuring notation}{
        name={goal structuring notation (GSN)},
        description={a graphical argumentation notation that explicitly represents the individual elements of any safety argument (requirements, claims, evidence and context) and the relationships that exist between these elements \cite{kelly2004goal}}
}

\newglossaryentry{known answer test}{
	name={known answer test (KAT)},
	description={a test of implementation correctness that employs pairs of parameter values and results that are calculated or otherwise known independently of the implementation being tested}
}

\newglossaryentry{model checking}{
	name={model checking},
	description={a form of automatic formal verification that employs a model checker tool to verify temporal behavior of a system
	as modeled by a symbolic state machine (transition system);
	a model checker may be used as an evidential tool to generate evidence items in an assurance case by O-ETB}
}

\newglossaryentry{O-ETB command (O-ETB)}{
	name={O-ETB command (O-ETB)},
	description={a command defined and accepted by the O-ETB command interpreter, either interactively or in a command procedure or script, to invoke an O-ETB function}
}

\newglossaryentry{Open Evidential Tool Bus}{
        name={Open Evidential Tool Bus (O-ETB)},
        description={an implementation of the ETB concept \cite{Rushby05} derived from AM-ETB \cite{CITADEL-D5.1} and adapted and extended for IoMT and MedSecurance by The Open Group \cite{O-ETB-UM}}
}

\newglossaryentry{privacy}{
	name={privacy},
	description={right of individuals to control or influence what information related to them may be collected and stored and by whom and to whom that information may be disclosed (source: ISO  TS 17574:2009)}
}

\newglossaryentry{protection profile}{
	name={protection profile (CC)},
	description={a requirements statement put together using CC constructs. Protection Profiles are generally published by governments for a specific technology type \cite{CC,ISOstd-15408-1}}
}

\newglossaryentry{reliability}{
	name={reliability},
	description={ability of a system or component to perform its required functions under stated conditions for a specified period of time (source: ISO/IEC 27040:2015)}
}

\newglossaryentry{REPOSITORY}{
        name={REPOSITORY (O-ETB)},
        description={a directory containing the persistent CASES Repository and the EVIDENCE Repository created and used by O-ETB \cite{O-ETB-UM}}
}

\newglossaryentry{resilience}{
	name={resilience},
	description={the condition of the system being able to avoid, absorb and/or manage dynamic adversarial conditions while completing assigned mission(s), and to reconstitute operational capabilities after casualties (source: IIC)}
}

\newglossaryentry{RESTful API}{
	name={RESTful API (or REST API)},
	description={an interface based on the REST architectural style that allows applications to communicate over the Internet}
}

\newglossaryentry{safety}{
	name={safety},
	description={the condition of the system operating without causing unacceptable risk of physical injury or damage to the health of people, either directly, or indirectly as a result of damage to property or to the environment (source: ISO/IEC Guide 55:1999)}
}

\newglossaryentry{security}{
	name={security},
	description={condition of the system being protected from unintended or unauthorized access, change or destruction. [Note: Security is a property of a system by which confidentiality, integrity, availability, accountability, authenticity, and reliability are achieved (ISO TR 15443-1:2012)}
}

\newglossaryentry{security target}{
	name={security target (CC)},
	description={implementation-dependent statement of security needs for a specific identified Target of Evaluation (TOE) \cite{CC,ISOstd-15408-1}}
}

\newglossaryentry{source repository}{
        name={source repository (O-ETB)},
        description={a directory structure containing the development directories of source code, documentation, build and test support for the O-ETB,
			and a template for the persistent database directories (KB, REPOSITORY and CAP) used and created by O-ETB \cite{O-ETB-UM}}
}

\newglossaryentry{stress test}{
	name={stress test},
	description={a test of a program or system that demands that it operate at high capacity, in volume and/or time, without failure}
}

\newglossaryentry{structured assurance case metamodel}{
	name={structured assurance case metamodel (SACM)},
	description={OMG standard metamodel for representing structured assurance cases and a graphical notation for depicting an assurance case \cite{SACM,SACM2}}                          
}

\newglossaryentry{system model}{
	name={system model (O-ETB)},
	description={a formal model of the architecture, properties and interactions of a system's components, used by O-ETB to provide information about the system that is needed during the assurance case argument pattern instantiation process}
}

\newglossaryentry{target of evaluation}{
	name={target of evaluation (CC)},
	description={an information system, part of a system or product, and all associated documentation, that is the subject of a security evaluation,
		in particular under the CC \cite{CC,ISOstd-15408-1}}
}

\newglossaryentry{tool agent}{
	name={tool agent (O-ETB)},
	description={a program invoked in a standard way by O-ETB that is intended to manage the execution of an external evidential tool,
		or performance of another assurance activity,
		to generate or validate an evidence item by invoking an external evidential tool or other assurance activity;
		the tool agent implements an interface to the external tool that is based on the specific data needs and interactions of that tool
		the tool agent updates the item in the EVIDENCE Repository that corresponds to an instance of the evidence caegory referenced by
		one or more particular locations in an assurance case \cite{O-ETB-UM}}
}

\newglossaryentry{trustworthiness}{
	name={trustworthiness},
	description={the degree of confidence one has that the system performs as expected with characteristics including safety, security, privacy, reliability and resilience in the face of environmental disruptions, human errors, system faults and attacks \cite{IISF}}
}

\newglossaryentry{workflow}{
	name={workflow},
	description={a sequence of industrial, administrative or other processes through which a piece of work passes from initiation to completion}
}

\newglossaryentry{workflow activity}{
	name={workflow activity},
	description={a specific task in a workflow system that is to be carried out by an automated or human agent; in the context of O-ETB, an assurance activity}
}

\newglossaryentry{workflow program}{
	name={workflow program},
	description={a sequence of statements expressed in a workflow language that defines the steps of a business process that are to be carried out as workflow activities by automated or human agents;
		O-ETB executes a workflow of assurance activities and tool invocations that is implicit in the procedure by which argument patterns are instantiated}
}

\newacronym{AC}{AC}{Assurance Case}
\newacronym{ACSL}{ACSL}{ANSI C Specification Language}
\newacronym{A-G}{A-G}{Assume-Guarantee}
\newacronym{API}{API}{Application Programming Interface}
\newacronym{CAE}{CAE}{Claims-Argument-Evidence}
\newacronym{CAP}{CAP}{Certification Assurance Package \cite{O-ETB-UM}}
\newacronym{CC}{CC}{Common Criteria for Information Technology Security Evaluation \cite{ISOstd-15408-1,ISOstd-15408-2,ISOstd-15408-3}}
\newacronym{CR}{CR}{CASES Repository (O-ETB) \cite{O-ETB-UM}}
\newacronym{CVE}{CVE}{Common Vulnerabilities Enumeration \cite{CVE}}
\newacronym{CWE}{CWE}{Common Weakness Enumeration \cite{CWE}}
\newacronym{CAPEC}{CAPEC}{Common Attack Pattern Enumeration and Classification \cite{CAPEC}}
\newacronym{CxC}{CxC}{Correct by Construction}
\newacronym{EAL}{EAL}{Evaluation Assurance Level \cite{ISOstd-15408-1}}
\newacronym{ER}{ER}{EVIDENCE Repository (O-ETB) \cite{O-ETB-UM}}
\newacronym{ETB}{ETB}{Evidential Tool Bus \cite{Rushby05}}
\newacronym{FV}{FV}{Formal Verification}
\newacronym{GSN}{GSN}{Goal Structuring Notation \cite{GSN_standard,GSN_standard_3}}
\newacronym{HiLARe}{HiLARe}{High-level ACSL Requirements}
\newacronym{IIRA}{IIRA}{Industrial Internet Reference Architecture \cite{IIRA}}
\newacronym{IISF}{IISF}{Industrial Internet Security Framework \cite{IISF}}
\newacronym{IIV}{IIV}{The Industrial Internet Vocabulary \cite{IIVocab}}
\newacronym{IoMT}{IoMT}{Internet of Medical Things \cite{MedSecurance-WWW}}
\newacronym{KB}{KB}{Knowledge Base \cite{O-ETB-UM}}
\newacronym{ltl}{LTL}{Linear-time Temporal Logic}
\newacronym{O-ETB}{O-ETB}{Open Evidential Tool Bus \cite{O-ETB-UM}}
\newacronym{OMG}{OMG}{Object Management Group}
\newacronym{MedSed}{MedSec}{MedSecurance Project}
\newacronym{MS}{MS}{MedSecurance Project}
\newacronym{PP}{PP}{Protection Profile \cite{CC,ISOstd-15408-1}}
\newacronym{REST}{REST}{Representational State Transfer}
\newacronym{RMV}{RMV}{Runtime Monitoring and Verification}
\newacronym{SACM}{SACM}{Structured Assurance Case Metamodel \cite{SACM,SACM2}}
\newacronym{SAR}{SAR}{Security Assurance Requirement \cite{CC,ISOstd-15408-1,ISOstd-15408-3}}
\newacronym{ST}{ST}{Security Target \cite{CC,ISOstd-15408-1}}
\newacronym{SFR}{SFR}{Security Functional Requirement \cite{CC,ISOstd-15408-1,ISOstd-15408-2}}
\newacronym{TA}{TA}{Tool Agent \cite{O-ETB-UM}}
\newacronym{TOE}{TOE}{Target of Evaluation \cite{CC,ISOstd-15408-1}}
\newacronym{TVRA}{TVRA}{Threat Vulnerability Risk Analysis}

%
% Don't forget to add e.g. ../../glossary-Assurance dependency to your Makefile if you
% are adding terms to a glossary while working on your document.

\newglossaryentry{analysis tool}{
	name={analysis tool},
	description={a tool that is used to perform some analysis or test the result of which may serve as evidence in support of an AC}
}

\newglossaryentry{argument pattern}{
	name={argument pattern},
	description={a parameterized template described in the GSN standard for an AC fragment possibly involving structural abstraction}
}

\newglossaryentry{assurance case}{
	name={assurance case},
	description={a documented body of evidence that provides a convincing and valid argument that a system is adequately dependable for a given application in a given environment; provides a set of auditable claims, arguments, and evidence created to support the claim that a defined system/service will satisfy particular requirements}
}

\newglossaryentry{assurance deficit}{
	name={assurance deficit},
	description={a documented body of evidence that provides a convincing and valid argument that a system is adequately dependable for a given application in a given environment; provides a set of auditable claims, arguments, and evidence created to support the claim that a defined system/service will satisfy particular requirements}
}

\newglossaryentry{assurance case argument pattern}{
        name={assurance case argument pattern},
        description={a parameterized expression representing logical relationships among goals (or claims) and evidence, expressed in a form of GSN with extensions for structural abstraction (see also: argument pattern}
}

\newglossaryentry{batch mode}{
        name={batch mode (O-ETB)},
        description={a mode of operation of O-ETB that is activated when the \texttt{-{}-command} or \texttt{-c} option is used to
			execute a single command or to run pre-defined command procedure or script (by executing the \texttt{proc} or \texttt{script}
			as the single command);
			other options that establish the context for the execution of the command are performed before the command;
			O-ETB exits after performing the specified command \cite{O-ETB-UM}}
}

\newglossaryentry{CAP artefact}{
	name={CAP artefact (O-ETB)},
	description={an object consisting of one or more files located in the CAP directory or a project directory within the CAP that has been produced to fulfill a requirement for certification evaluation}
}

\newglossaryentry{CAP Renderer}{
	name={CAP Renderer (O-ETB)},
	description={a component of O-ETB that renders information in the CASES Repository and the EVIDENCE Repository in an external form suitable for review as part of a certification evaluation}
}

\newglossaryentry{CASES Repository}{
	name={CASES Repository (O-ETB)},
	description={a persistent store created by O-ETB that contains modular ACs instantiated from one or more argument patterns and referring to evidence records in the ER \cite{O-ETB-UM}}
}

\newglossaryentry{certification assurance package}{
        name={certification assurance package (CAP) (O-ETB)},
        description={a collection of documents (textual and/or graphical and electronically searchable),
			including those generated by the O-ETB,
			required for evaluation that present required evidence in a form acceptable to a certification authority \cite{O-ETB-UM}}
}

\newglossaryentry{certification project}{
        name={certification project},
        description={a subdivision within the CAP for sets of artifacts corresponding to different certification activities \cite{O-ETB-UM}}
}

\newglossaryentry{certification requirements}{
	name={certification requirements},
	description={a set of requirements prescribing methods and/or the content and presentation of evidence needed to justify granting certification according to some particular domain or standard}
}

\newglossaryentry{claims-argument-evidence}{
        name={claims-argument-evidence (CAE)},
        description={a structured approach which aids the visibility, review and assessment of safety case documentation \cite{CAE}}
}

\newglossaryentry{command mode}{
        name={command mode (O-ETB)},
        description={a command-based interface that can be used interactively and/or in a batch mode to execute single commands or to run pre-defined command scripts \cite{O-ETB-UM}}
}

\newglossaryentry{command procedure}{
        name={command procedure (O-ETB)},
        description={a pre-defined named sequence of O-ETB commands, defined along with other such procedures
			in the \texttt{procs\_etb.pl} file,
			that can be invoked by using the \texttt{proc} command either in the interactive command interpreter
			or in a \texttt{-{}-command} option in batch mode \cite{O-ETB-UM}}
}

\newglossaryentry{command script}{
        name={command script (O-ETB)},
        description={a user-defined sequence of O-ETB commands contained in a named file,
			that can be invoked by using the \texttt{script} command either in the interactive command interpreter
			or in a \texttt{-{}-command} option in batch mode \cite{O-ETB-UM}}
}

\newglossaryentry{common attack pattern enumeration and classification}{
        name={common attack pattern enumeration and classification (CAPEC)},
        description={a comprehensive dictionary of known patterns of attack employed by adversaries to exploit known weaknesses in cyber-enabled capabilities \cite{CAPEC}}
}

\newglossaryentry{common criteria}{
        name={common criteria (CC)},
        description={Common Criteria for Information Security Technology Evaluation \cite{ISOstd-15408-1,ISOstd-15408-2,ISOstd-15408-3}}
}

\newglossaryentry{common vulnerabilities enumeration}{
        name={common vulnerabilities enumeration (CVE)},
        description={catalog of publicly disclosed cybersecurity vulnerabilities \cite{CVE}}
}

\newglossaryentry{common weaknesses enumeration}{
        name={common weaknesses enumeration (CWE)},
        description={the most severe and prevalent weaknesses behind the CVE records \cite{CWE}}
}

\newglossaryentry{correct by construction}{
	name={correct by construction (CxC)},
	description={rigorous techniques applied to the construction of some artifact or product that logically ensure that the item correctly adheres to a specified property or properties}
}

\newglossaryentry{deductive verification}{
	name={deductive verification},
	description={deductive software verification aims at formally verifying that all possible behaviors of a given program satisfy formally defined, possibly complex properties, where the verification process is based on some form of logical inference, i.e., deduction}
}

\newglossaryentry{ETB Core Workflow Engine}{
	name={ETB Core Workflow Engine (O-ETB)},
	description={a component of O-ETB that is the top-level driver for the instantiation of assurance case argument patterns and the passing of actual parameter values \cite{O-ETB-UM}}
}

\newglossaryentry{ETB Developer/Administrator}{
	name={ETB Developer/Administrator (O-ETB)},
	description={the role of a person that develops ETB implementation or administers an ETB implementation}
}

\newglossaryentry{evaluation assurance level}{
        name={evaluation assurance level (EAL) (CC)},
        description={set of assurance requirements that represent a point on the Common Criteria predefined assurance scale (EALs 1..7) \cite{CC,ISOstd-15408-1}}
}

\newglossaryentry{evaluator (certifier)}{
	name={evaluator (certifier)},
	description={a qualified expert who examines evidence presented in support of an application by a developer to certify a product}
}

\newglossaryentry{evidence}{
	name={evidence},
	description={concrete empirical data that may be used to support arguments for claims}
}

\newglossaryentry{evidence category}{
	name={evidence category (O-ETB)},
	description={an identifier for a set of values that a particular O-ETB evidence item may represent and which is produced by a particular method or tool \cite{O-ETB-UM}}
}

\newglossaryentry{EVIDENCE Repository}{
	name={EVIDENCE Repository (O-ETB)},
	description={a persistent store created by O-ETB and updated by tool agents that contains evidence records that represent placeholders for evidence or the concrete result of an assurance activity or the output of an evidential tool \cite{O-ETB-UM}}
}

\newglossaryentry{evidence specifier}{
	name={evidence specifier (O-ETB)},
	description={unique identifier for an instance of an evidence category occurring in the EVIDENCE Repository, see D2.2 p26 \cite{O-ETB-UM}}
}

\newglossaryentry{evidential tool bus}{
	name={evidential tool bus (ETB)},
	description={a system such as O-ETB that constructs an assurance case and supporting evidence according to a collection of argument patterns and invokes and records the results from external evidential tools \cite{Rushby05}}
}

\newglossaryentry{Frama-C}{
	name={Frama-C},
	description={Framework for Modular Analysis of C programs;
		an open-source extensible and collaborative platform dedicated to source-code analysis of C software}
}

\newglossaryentry{Frama-PLC}{
	name={Frama-PLC},
	description={performs syntactical and semantic checks of PLC application software}
}

\newglossaryentry{goal structuring notation}{
        name={goal structuring notation (GSN)},
        description={a graphical argumentation notation that explicitly represents the individual elements of any safety argument (requirements, claims, evidence and context) and the relationships that exist between these elements \cite{kelly2004goal}}
}

\newglossaryentry{known answer test}{
	name={known answer test (KAT)},
	description={a test of implementation correctness that employs pairs of parameter values and results that are calculated or otherwise known independently of the implementation being tested}
}

\newglossaryentry{model checking}{
	name={model checking},
	description={a form of automatic formal verification that employs a model checker tool to verify temporal behavior of a system
	as modeled by a symbolic state machine (transition system);
	a model checker may be used as an evidential tool to generate evidence items in an assurance case by O-ETB}
}

\newglossaryentry{O-ETB command (O-ETB)}{
	name={O-ETB command (O-ETB)},
	description={a command defined and accepted by the O-ETB command interpreter, either interactively or in a command procedure or script, to invoke an O-ETB function}
}

\newglossaryentry{Open Evidential Tool Bus}{
        name={Open Evidential Tool Bus (O-ETB)},
        description={an implementation of the ETB concept \cite{Rushby05} derived from AM-ETB \cite{CITADEL-D5.1} and adapted and extended for IoMT and MedSecurance by The Open Group \cite{O-ETB-UM}}
}

\newglossaryentry{privacy}{
	name={privacy},
	description={right of individuals to control or influence what information related to them may be collected and stored and by whom and to whom that information may be disclosed (source: ISO  TS 17574:2009)}
}

\newglossaryentry{protection profile}{
	name={protection profile (CC)},
	description={a requirements statement put together using CC constructs. Protection Profiles are generally published by governments for a specific technology type \cite{CC,ISOstd-15408-1}}
}

\newglossaryentry{reliability}{
	name={reliability},
	description={ability of a system or component to perform its required functions under stated conditions for a specified period of time (source: ISO/IEC 27040:2015)}
}

\newglossaryentry{REPOSITORY}{
        name={REPOSITORY (O-ETB)},
        description={a directory containing the persistent CASES Repository and the EVIDENCE Repository created and used by O-ETB \cite{O-ETB-UM}}
}

\newglossaryentry{resilience}{
	name={resilience},
	description={the condition of the system being able to avoid, absorb and/or manage dynamic adversarial conditions while completing assigned mission(s), and to reconstitute operational capabilities after casualties (source: IIC)}
}

\newglossaryentry{RESTful API}{
	name={RESTful API (or REST API)},
	description={an interface based on the REST architectural style that allows applications to communicate over the Internet}
}

\newglossaryentry{safety}{
	name={safety},
	description={the condition of the system operating without causing unacceptable risk of physical injury or damage to the health of people, either directly, or indirectly as a result of damage to property or to the environment (source: ISO/IEC Guide 55:1999)}
}

\newglossaryentry{security}{
	name={security},
	description={condition of the system being protected from unintended or unauthorized access, change or destruction. [Note: Security is a property of a system by which confidentiality, integrity, availability, accountability, authenticity, and reliability are achieved (ISO TR 15443-1:2012)}
}

\newglossaryentry{security target}{
	name={security target (CC)},
	description={implementation-dependent statement of security needs for a specific identified Target of Evaluation (TOE) \cite{CC,ISOstd-15408-1}}
}

\newglossaryentry{source repository}{
        name={source repository (O-ETB)},
        description={a directory structure containing the development directories of source code, documentation, build and test support for the O-ETB,
			and a template for the persistent database directories (KB, REPOSITORY and CAP) used and created by O-ETB \cite{O-ETB-UM}}
}

\newglossaryentry{stress test}{
	name={stress test},
	description={a test of a program or system that demands that it operate at high capacity, in volume and/or time, without failure}
}

\newglossaryentry{structured assurance case metamodel}{
	name={structured assurance case metamodel (SACM)},
	description={OMG standard metamodel for representing structured assurance cases and a graphical notation for depicting an assurance case \cite{SACM,SACM2}}                          
}

\newglossaryentry{system model}{
	name={system model (O-ETB)},
	description={a formal model of the architecture, properties and interactions of a system's components, used by O-ETB to provide information about the system that is needed during the assurance case argument pattern instantiation process}
}

\newglossaryentry{target of evaluation}{
	name={target of evaluation (CC)},
	description={an information system, part of a system or product, and all associated documentation, that is the subject of a security evaluation,
		in particular under the CC \cite{CC,ISOstd-15408-1}}
}

\newglossaryentry{tool agent}{
	name={tool agent (O-ETB)},
	description={a program invoked in a standard way by O-ETB that is intended to manage the execution of an external evidential tool,
		or performance of another assurance activity,
		to generate or validate an evidence item by invoking an external evidential tool or other assurance activity;
		the tool agent implements an interface to the external tool that is based on the specific data needs and interactions of that tool
		the tool agent updates the item in the EVIDENCE Repository that corresponds to an instance of the evidence caegory referenced by
		one or more particular locations in an assurance case \cite{O-ETB-UM}}
}

\newglossaryentry{trustworthiness}{
	name={trustworthiness},
	description={the degree of confidence one has that the system performs as expected with characteristics including safety, security, privacy, reliability and resilience in the face of environmental disruptions, human errors, system faults and attacks \cite{IISF}}
}

\newglossaryentry{workflow}{
	name={workflow},
	description={a sequence of industrial, administrative or other processes through which a piece of work passes from initiation to completion}
}

\newglossaryentry{workflow activity}{
	name={workflow activity},
	description={a specific task in a workflow system that is to be carried out by an automated or human agent; in the context of O-ETB, an assurance activity}
}

\newglossaryentry{workflow program}{
	name={workflow program},
	description={a sequence of statements expressed in a workflow language that defines the steps of a business process that are to be carried out as workflow activities by automated or human agents;
		O-ETB executes a workflow of assurance activities and tool invocations that is implicit in the procedure by which argument patterns are instantiated}
}

\newacronym{AC}{AC}{Assurance Case}
\newacronym{ACSL}{ACSL}{ANSI C Specification Language}
\newacronym{A-G}{A-G}{Assume-Guarantee}
\newacronym{API}{API}{Application Programming Interface}
\newacronym{CAE}{CAE}{Claims-Argument-Evidence}
\newacronym{CAP}{CAP}{Certification Assurance Package \cite{O-ETB-UM}}
\newacronym{CC}{CC}{Common Criteria for Information Technology Security Evaluation \cite{ISOstd-15408-1,ISOstd-15408-2,ISOstd-15408-3}}
\newacronym{CR}{CR}{CASES Repository (O-ETB) \cite{O-ETB-UM}}
\newacronym{CVE}{CVE}{Common Vulnerabilities Enumeration \cite{CVE}}
\newacronym{CWE}{CWE}{Common Weakness Enumeration \cite{CWE}}
\newacronym{CAPEC}{CAPEC}{Common Attack Pattern Enumeration and Classification \cite{CAPEC}}
\newacronym{CxC}{CxC}{Correct by Construction}
\newacronym{EAL}{EAL}{Evaluation Assurance Level \cite{ISOstd-15408-1}}
\newacronym{ER}{ER}{EVIDENCE Repository (O-ETB) \cite{O-ETB-UM}}
\newacronym{ETB}{ETB}{Evidential Tool Bus \cite{Rushby05}}
\newacronym{FV}{FV}{Formal Verification}
\newacronym{GSN}{GSN}{Goal Structuring Notation \cite{GSN_standard,GSN_standard_3}}
\newacronym{HiLARe}{HiLARe}{High-level ACSL Requirements}
\newacronym{IIRA}{IIRA}{Industrial Internet Reference Architecture \cite{IIRA}}
\newacronym{IISF}{IISF}{Industrial Internet Security Framework \cite{IISF}}
\newacronym{IIV}{IIV}{The Industrial Internet Vocabulary \cite{IIVocab}}
\newacronym{IoMT}{IoMT}{Internet of Medical Things \cite{MedSecurance-WWW}}
\newacronym{KB}{KB}{Knowledge Base \cite{O-ETB-UM}}
\newacronym{ltl}{LTL}{Linear-time Temporal Logic}
\newacronym{O-ETB}{O-ETB}{Open Evidential Tool Bus \cite{O-ETB-UM}}
\newacronym{OMG}{OMG}{Object Management Group}
\newacronym{MedSed}{MedSec}{MedSecurance Project}
\newacronym{MS}{MS}{MedSecurance Project}
\newacronym{PP}{PP}{Protection Profile \cite{CC,ISOstd-15408-1}}
\newacronym{REST}{REST}{Representational State Transfer}
\newacronym{RMV}{RMV}{Runtime Monitoring and Verification}
\newacronym{SACM}{SACM}{Structured Assurance Case Metamodel \cite{SACM,SACM2}}
\newacronym{SAR}{SAR}{Security Assurance Requirement \cite{CC,ISOstd-15408-1,ISOstd-15408-3}}
\newacronym{ST}{ST}{Security Target \cite{CC,ISOstd-15408-1}}
\newacronym{SFR}{SFR}{Security Functional Requirement \cite{CC,ISOstd-15408-1,ISOstd-15408-2}}
\newacronym{TA}{TA}{Tool Agent \cite{O-ETB-UM}}
\newacronym{TOE}{TOE}{Target of Evaluation \cite{CC,ISOstd-15408-1}}
\newacronym{TVRA}{TVRA}{Threat Vulnerability Risk Analysis}

%
% Don't forget to add e.g. ../../glossary-Assurance dependency to your Makefile if you
% are adding terms to a glossary while working on your document.

\newglossaryentry{analysis tool}{
	name={analysis tool},
	description={a tool that is used to perform some analysis or test the result of which may serve as evidence in support of an AC}
}

\newglossaryentry{argument pattern}{
	name={argument pattern},
	description={a parameterized template described in the GSN standard for an AC fragment possibly involving structural abstraction}
}

\newglossaryentry{assurance case}{
	name={assurance case},
	description={a documented body of evidence that provides a convincing and valid argument that a system is adequately dependable for a given application in a given environment; provides a set of auditable claims, arguments, and evidence created to support the claim that a defined system/service will satisfy particular requirements}
}

\newglossaryentry{assurance deficit}{
	name={assurance deficit},
	description={a documented body of evidence that provides a convincing and valid argument that a system is adequately dependable for a given application in a given environment; provides a set of auditable claims, arguments, and evidence created to support the claim that a defined system/service will satisfy particular requirements}
}

\newglossaryentry{assurance case argument pattern}{
        name={assurance case argument pattern},
        description={a parameterized expression representing logical relationships among goals (or claims) and evidence, expressed in a form of GSN with extensions for structural abstraction (see also: argument pattern}
}

\newglossaryentry{batch mode}{
        name={batch mode (O-ETB)},
        description={a mode of operation of O-ETB that is activated when the \texttt{-{}-command} or \texttt{-c} option is used to
			execute a single command or to run pre-defined command procedure or script (by executing the \texttt{proc} or \texttt{script}
			as the single command);
			other options that establish the context for the execution of the command are performed before the command;
			O-ETB exits after performing the specified command \cite{O-ETB-UM}}
}

\newglossaryentry{CAP artefact}{
	name={CAP artefact (O-ETB)},
	description={an object consisting of one or more files located in the CAP directory or a project directory within the CAP that has been produced to fulfill a requirement for certification evaluation}
}

\newglossaryentry{CAP Renderer}{
	name={CAP Renderer (O-ETB)},
	description={a component of O-ETB that renders information in the CASES Repository and the EVIDENCE Repository in an external form suitable for review as part of a certification evaluation}
}

\newglossaryentry{CASES Repository}{
	name={CASES Repository (O-ETB)},
	description={a persistent store created by O-ETB that contains modular ACs instantiated from one or more argument patterns and referring to evidence records in the ER \cite{O-ETB-UM}}
}

\newglossaryentry{certification assurance package}{
        name={certification assurance package (CAP) (O-ETB)},
        description={a collection of documents (textual and/or graphical and electronically searchable),
			including those generated by the O-ETB,
			required for evaluation that present required evidence in a form acceptable to a certification authority \cite{O-ETB-UM}}
}

\newglossaryentry{certification project}{
        name={certification project},
        description={a subdivision within the CAP for sets of artifacts corresponding to different certification activities \cite{O-ETB-UM}}
}

\newglossaryentry{certification requirements}{
	name={certification requirements},
	description={a set of requirements prescribing methods and/or the content and presentation of evidence needed to justify granting certification according to some particular domain or standard}
}

\newglossaryentry{claims-argument-evidence}{
        name={claims-argument-evidence (CAE)},
        description={a structured approach which aids the visibility, review and assessment of safety case documentation \cite{CAE}}
}

\newglossaryentry{command mode}{
        name={command mode (O-ETB)},
        description={a command-based interface that can be used interactively and/or in a batch mode to execute single commands or to run pre-defined command scripts \cite{O-ETB-UM}}
}

\newglossaryentry{command procedure}{
        name={command procedure (O-ETB)},
        description={a pre-defined named sequence of O-ETB commands, defined along with other such procedures
			in the \texttt{procs\_etb.pl} file,
			that can be invoked by using the \texttt{proc} command either in the interactive command interpreter
			or in a \texttt{-{}-command} option in batch mode \cite{O-ETB-UM}}
}

\newglossaryentry{command script}{
        name={command script (O-ETB)},
        description={a user-defined sequence of O-ETB commands contained in a named file,
			that can be invoked by using the \texttt{script} command either in the interactive command interpreter
			or in a \texttt{-{}-command} option in batch mode \cite{O-ETB-UM}}
}

\newglossaryentry{common attack pattern enumeration and classification}{
        name={common attack pattern enumeration and classification (CAPEC)},
        description={a comprehensive dictionary of known patterns of attack employed by adversaries to exploit known weaknesses in cyber-enabled capabilities \cite{CAPEC}}
}

\newglossaryentry{common criteria}{
        name={common criteria (CC)},
        description={Common Criteria for Information Security Technology Evaluation \cite{ISOstd-15408-1,ISOstd-15408-2,ISOstd-15408-3}}
}

\newglossaryentry{common vulnerabilities enumeration}{
        name={common vulnerabilities enumeration (CVE)},
        description={catalog of publicly disclosed cybersecurity vulnerabilities \cite{CVE}}
}

\newglossaryentry{common weaknesses enumeration}{
        name={common weaknesses enumeration (CWE)},
        description={the most severe and prevalent weaknesses behind the CVE records \cite{CWE}}
}

\newglossaryentry{correct by construction}{
	name={correct by construction (CxC)},
	description={rigorous techniques applied to the construction of some artifact or product that logically ensure that the item correctly adheres to a specified property or properties}
}

\newglossaryentry{deductive verification}{
	name={deductive verification},
	description={deductive software verification aims at formally verifying that all possible behaviors of a given program satisfy formally defined, possibly complex properties, where the verification process is based on some form of logical inference, i.e., deduction}
}

\newglossaryentry{ETB Core Workflow Engine}{
	name={ETB Core Workflow Engine (O-ETB)},
	description={a component of O-ETB that is the top-level driver for the instantiation of assurance case argument patterns and the passing of actual parameter values \cite{O-ETB-UM}}
}

\newglossaryentry{ETB Developer/Administrator}{
	name={ETB Developer/Administrator (O-ETB)},
	description={the role of a person that develops ETB implementation or administers an ETB implementation}
}

\newglossaryentry{evaluation assurance level}{
        name={evaluation assurance level (EAL) (CC)},
        description={set of assurance requirements that represent a point on the Common Criteria predefined assurance scale (EALs 1..7) \cite{CC,ISOstd-15408-1}}
}

\newglossaryentry{evaluator (certifier)}{
	name={evaluator (certifier)},
	description={a qualified expert who examines evidence presented in support of an application by a developer to certify a product}
}

\newglossaryentry{evidence}{
	name={evidence},
	description={concrete empirical data that may be used to support arguments for claims}
}

\newglossaryentry{evidence category}{
	name={evidence category (O-ETB)},
	description={an identifier for a set of values that a particular O-ETB evidence item may represent and which is produced by a particular method or tool \cite{O-ETB-UM}}
}

\newglossaryentry{EVIDENCE Repository}{
	name={EVIDENCE Repository (O-ETB)},
	description={a persistent store created by O-ETB and updated by tool agents that contains evidence records that represent placeholders for evidence or the concrete result of an assurance activity or the output of an evidential tool \cite{O-ETB-UM}}
}

\newglossaryentry{evidence specifier}{
	name={evidence specifier (O-ETB)},
	description={unique identifier for an instance of an evidence category occurring in the EVIDENCE Repository, see D2.2 p26 \cite{O-ETB-UM}}
}

\newglossaryentry{evidential tool bus}{
	name={evidential tool bus (ETB)},
	description={a system such as O-ETB that constructs an assurance case and supporting evidence according to a collection of argument patterns and invokes and records the results from external evidential tools \cite{Rushby05}}
}

\newglossaryentry{Frama-C}{
	name={Frama-C},
	description={Framework for Modular Analysis of C programs;
		an open-source extensible and collaborative platform dedicated to source-code analysis of C software}
}

\newglossaryentry{Frama-PLC}{
	name={Frama-PLC},
	description={performs syntactical and semantic checks of PLC application software}
}

\newglossaryentry{goal structuring notation}{
        name={goal structuring notation (GSN)},
        description={a graphical argumentation notation that explicitly represents the individual elements of any safety argument (requirements, claims, evidence and context) and the relationships that exist between these elements \cite{kelly2004goal}}
}

\newglossaryentry{known answer test}{
	name={known answer test (KAT)},
	description={a test of implementation correctness that employs pairs of parameter values and results that are calculated or otherwise known independently of the implementation being tested}
}

\newglossaryentry{model checking}{
	name={model checking},
	description={a form of automatic formal verification that employs a model checker tool to verify temporal behavior of a system
	as modeled by a symbolic state machine (transition system);
	a model checker may be used as an evidential tool to generate evidence items in an assurance case by O-ETB}
}

\newglossaryentry{O-ETB command (O-ETB)}{
	name={O-ETB command (O-ETB)},
	description={a command defined and accepted by the O-ETB command interpreter, either interactively or in a command procedure or script, to invoke an O-ETB function}
}

\newglossaryentry{Open Evidential Tool Bus}{
        name={Open Evidential Tool Bus (O-ETB)},
        description={an implementation of the ETB concept \cite{Rushby05} derived from AM-ETB \cite{CITADEL-D5.1} and adapted and extended for IoMT and MedSecurance by The Open Group \cite{O-ETB-UM}}
}

\newglossaryentry{privacy}{
	name={privacy},
	description={right of individuals to control or influence what information related to them may be collected and stored and by whom and to whom that information may be disclosed (source: ISO  TS 17574:2009)}
}

\newglossaryentry{protection profile}{
	name={protection profile (CC)},
	description={a requirements statement put together using CC constructs. Protection Profiles are generally published by governments for a specific technology type \cite{CC,ISOstd-15408-1}}
}

\newglossaryentry{reliability}{
	name={reliability},
	description={ability of a system or component to perform its required functions under stated conditions for a specified period of time (source: ISO/IEC 27040:2015)}
}

\newglossaryentry{REPOSITORY}{
        name={REPOSITORY (O-ETB)},
        description={a directory containing the persistent CASES Repository and the EVIDENCE Repository created and used by O-ETB \cite{O-ETB-UM}}
}

\newglossaryentry{resilience}{
	name={resilience},
	description={the condition of the system being able to avoid, absorb and/or manage dynamic adversarial conditions while completing assigned mission(s), and to reconstitute operational capabilities after casualties (source: IIC)}
}

\newglossaryentry{RESTful API}{
	name={RESTful API (or REST API)},
	description={an interface based on the REST architectural style that allows applications to communicate over the Internet}
}

\newglossaryentry{safety}{
	name={safety},
	description={the condition of the system operating without causing unacceptable risk of physical injury or damage to the health of people, either directly, or indirectly as a result of damage to property or to the environment (source: ISO/IEC Guide 55:1999)}
}

\newglossaryentry{security}{
	name={security},
	description={condition of the system being protected from unintended or unauthorized access, change or destruction. [Note: Security is a property of a system by which confidentiality, integrity, availability, accountability, authenticity, and reliability are achieved (ISO TR 15443-1:2012)}
}

\newglossaryentry{security target}{
	name={security target (CC)},
	description={implementation-dependent statement of security needs for a specific identified Target of Evaluation (TOE) \cite{CC,ISOstd-15408-1}}
}

\newglossaryentry{source repository}{
        name={source repository (O-ETB)},
        description={a directory structure containing the development directories of source code, documentation, build and test support for the O-ETB,
			and a template for the persistent database directories (KB, REPOSITORY and CAP) used and created by O-ETB \cite{O-ETB-UM}}
}

\newglossaryentry{stress test}{
	name={stress test},
	description={a test of a program or system that demands that it operate at high capacity, in volume and/or time, without failure}
}

\newglossaryentry{structured assurance case metamodel}{
	name={structured assurance case metamodel (SACM)},
	description={OMG standard metamodel for representing structured assurance cases and a graphical notation for depicting an assurance case \cite{SACM,SACM2}}                          
}

\newglossaryentry{system model}{
	name={system model (O-ETB)},
	description={a formal model of the architecture, properties and interactions of a system's components, used by O-ETB to provide information about the system that is needed during the assurance case argument pattern instantiation process}
}

\newglossaryentry{target of evaluation}{
	name={target of evaluation (CC)},
	description={an information system, part of a system or product, and all associated documentation, that is the subject of a security evaluation,
		in particular under the CC \cite{CC,ISOstd-15408-1}}
}

\newglossaryentry{tool agent}{
	name={tool agent (O-ETB)},
	description={a program invoked in a standard way by O-ETB that is intended to manage the execution of an external evidential tool,
		or performance of another assurance activity,
		to generate or validate an evidence item by invoking an external evidential tool or other assurance activity;
		the tool agent implements an interface to the external tool that is based on the specific data needs and interactions of that tool
		the tool agent updates the item in the EVIDENCE Repository that corresponds to an instance of the evidence caegory referenced by
		one or more particular locations in an assurance case \cite{O-ETB-UM}}
}

\newglossaryentry{trustworthiness}{
	name={trustworthiness},
	description={the degree of confidence one has that the system performs as expected with characteristics including safety, security, privacy, reliability and resilience in the face of environmental disruptions, human errors, system faults and attacks \cite{IISF}}
}

\newglossaryentry{workflow}{
	name={workflow},
	description={a sequence of industrial, administrative or other processes through which a piece of work passes from initiation to completion}
}

\newglossaryentry{workflow activity}{
	name={workflow activity},
	description={a specific task in a workflow system that is to be carried out by an automated or human agent; in the context of O-ETB, an assurance activity}
}

\newglossaryentry{workflow program}{
	name={workflow program},
	description={a sequence of statements expressed in a workflow language that defines the steps of a business process that are to be carried out as workflow activities by automated or human agents;
		O-ETB executes a workflow of assurance activities and tool invocations that is implicit in the procedure by which argument patterns are instantiated}
}

\newacronym{AC}{AC}{Assurance Case}
\newacronym{ACSL}{ACSL}{ANSI C Specification Language}
\newacronym{A-G}{A-G}{Assume-Guarantee}
\newacronym{API}{API}{Application Programming Interface}
\newacronym{CAE}{CAE}{Claims-Argument-Evidence}
\newacronym{CAP}{CAP}{Certification Assurance Package \cite{O-ETB-UM}}
\newacronym{CC}{CC}{Common Criteria for Information Technology Security Evaluation \cite{ISOstd-15408-1,ISOstd-15408-2,ISOstd-15408-3}}
\newacronym{CR}{CR}{CASES Repository (O-ETB) \cite{O-ETB-UM}}
\newacronym{CVE}{CVE}{Common Vulnerabilities Enumeration \cite{CVE}}
\newacronym{CWE}{CWE}{Common Weakness Enumeration \cite{CWE}}
\newacronym{CAPEC}{CAPEC}{Common Attack Pattern Enumeration and Classification \cite{CAPEC}}
\newacronym{CxC}{CxC}{Correct by Construction}
\newacronym{EAL}{EAL}{Evaluation Assurance Level \cite{ISOstd-15408-1}}
\newacronym{ER}{ER}{EVIDENCE Repository (O-ETB) \cite{O-ETB-UM}}
\newacronym{ETB}{ETB}{Evidential Tool Bus \cite{Rushby05}}
\newacronym{FV}{FV}{Formal Verification}
\newacronym{GSN}{GSN}{Goal Structuring Notation \cite{GSN_standard,GSN_standard_3}}
\newacronym{HiLARe}{HiLARe}{High-level ACSL Requirements}
\newacronym{IIRA}{IIRA}{Industrial Internet Reference Architecture \cite{IIRA}}
\newacronym{IISF}{IISF}{Industrial Internet Security Framework \cite{IISF}}
\newacronym{IIV}{IIV}{The Industrial Internet Vocabulary \cite{IIVocab}}
\newacronym{IoMT}{IoMT}{Internet of Medical Things \cite{MedSecurance-WWW}}
\newacronym{KB}{KB}{Knowledge Base \cite{O-ETB-UM}}
\newacronym{ltl}{LTL}{Linear-time Temporal Logic}
\newacronym{O-ETB}{O-ETB}{Open Evidential Tool Bus \cite{O-ETB-UM}}
\newacronym{OMG}{OMG}{Object Management Group}
\newacronym{MedSed}{MedSec}{MedSecurance Project}
\newacronym{MS}{MS}{MedSecurance Project}
\newacronym{PP}{PP}{Protection Profile \cite{CC,ISOstd-15408-1}}
\newacronym{REST}{REST}{Representational State Transfer}
\newacronym{RMV}{RMV}{Runtime Monitoring and Verification}
\newacronym{SACM}{SACM}{Structured Assurance Case Metamodel \cite{SACM,SACM2}}
\newacronym{SAR}{SAR}{Security Assurance Requirement \cite{CC,ISOstd-15408-1,ISOstd-15408-3}}
\newacronym{ST}{ST}{Security Target \cite{CC,ISOstd-15408-1}}
\newacronym{SFR}{SFR}{Security Functional Requirement \cite{CC,ISOstd-15408-1,ISOstd-15408-2}}
\newacronym{TA}{TA}{Tool Agent \cite{O-ETB-UM}}
\newacronym{TOE}{TOE}{Target of Evaluation \cite{CC,ISOstd-15408-1}}
\newacronym{TVRA}{TVRA}{Threat Vulnerability Risk Analysis}
